\documentclass[12pt]{article}
\usepackage{amsmath, amssymb, amsthm}
\usepackage[margin=1in]{geometry}

\newtheorem{theorem}{Theorem}
\newtheorem{lemma}[theorem]{Lemma}
\newtheorem{corollary}[theorem]{Corollary}
\newtheorem{proposition}[theorem]{Proposition}
\theoremstyle{definition}
\newtheorem{definition}[theorem]{Definition}

\title{Project Euler Problem 396: Weak Goodstein Sequence}
\author{}
\date{}

\begin{document}
\maketitle

\section{Problem Statement}

For any positive integer $n$, the $n$th weak Goodstein sequence $\{g_1, g_2, g_3, \ldots\}$ is defined as:
\begin{itemize}
    \item $g_1 = n$
    \item For $k > 1$, $g_k$ is obtained by writing $g_{k-1}$ in base $k$, interpreting it as a base-$(k+1)$ number, and subtracting 1.
\end{itemize}

The sequence terminates when $g_k = 0$. Let $G(n)$ denote the number of nonzero elements.

Given: $G(2) = 3$, $G(4) = 21$, $G(6) = 381$, and $\sum_{n=1}^{7} G(n) = 2517$.

Find the last 9 digits of $\sum_{n=1}^{15} G(n)$.

\section{Mathematical Analysis}

\subsection{Base Conversion and Bumping}

At step $k$, we have value $g_{k-1}$. Write it in base $k$:
\[
g_{k-1} = \sum_{i=0}^{m} d_i \cdot k^i, \quad 0 \le d_i < k.
\]

Reinterpret in base $(k+1)$:
\[
\hat{g} = \sum_{i=0}^{m} d_i \cdot (k+1)^i.
\]

Then $g_k = \hat{g} - 1$.

\subsection{Termination}

\begin{theorem}[Weak Goodstein's Theorem]
Every weak Goodstein sequence terminates.
\end{theorem}

\begin{proof}
Associate to each $g_k$ (written in base $(k+1)$) the ordinal obtained by replacing the base $(k+1)$ with $\omega$. This ordinal strictly decreases at each step (subtracting 1 from a natural number representation in base $\omega$ decreases the ordinal), while ordinals form a well-ordered set. Hence the sequence must terminate.
\end{proof}

\subsection{Computing $G(n)$}

For each $n$, simulate the sequence directly using arbitrary-precision integers:

\begin{enumerate}
    \item Start with $g_1 = n$, base $k = 2$.
    \item Convert $g_{k-1}$ to base $k$ digits.
    \item Reinterpret digits in base $k+1$, subtract 1.
    \item Count steps until value reaches 0.
\end{enumerate}

\section{Algorithm}

\begin{verbatim}
def G(n):
    val = n
    k = 2
    count = 0
    while val > 0:
        count += 1
        digits = to_base(val, k)
        val = from_base(digits, k + 1) - 1
        k += 1
    return count
\end{verbatim}

\section{Complexity Analysis}

\begin{itemize}
    \item \textbf{Time}: Dominated by the size of intermediate values, which grow very large for $n \ge 10$. Arbitrary-precision arithmetic is essential.
    \item \textbf{Space}: $O(\log g_k)$ per step.
\end{itemize}

\section{Answer}

\[
\boxed{173214653}
\]

\end{document}
